\newtoggle{withbonus} \toggletrue{withbonus} % set true to show bonus points in grade table

% setting underlying course name (Grundlagenfach Mathematik, §-Unterricht, \ldots)
\newcommand{\mycourse}{Ergänzungsfach Informatik}
% name of the class
\newcommand{\myclass}{4. Klassen}
% set date
\newcommand{\mydate}{06. November 2025}

\title{\textbf{Kolmogorov-Komplexität}}
\author{Thomas Graf\\
	\mycourse, \myclass}
\date{\mydate}
\maketitle
\thispagestyle{headandfoot}

\begin{center}
	\fbox{\parbox{140mm}{\centering
			\begin{itemize}
				\item Dauer der Prüfung: $60$ Minuten
				      %\item \textbf{Zeige in jeder Aufgabe unbedingt Deine Schritte!}
				\item Erlaubte Hilfsmittel:
				      \begin{itemize}
					      \item Schreibzeug
					            % \item Taschenrechner
				      \end{itemize}
			\end{itemize}}}
\end{center}
\vspace{10mm}

\begin{center}
	Vorname: \rule{45mm}{0.8pt}\qquad Nachname: \rule{45mm}{0.8pt}
\end{center}
\vspace{20mm}

\begin{center}
	\iftoggle{withbonus} {
		\combinedgradetable[h][questions]
	} {
		\gradetable[h][questions]
	}
\end{center}
\vspace{15mm}
\begin{center}
	Note:
	\vspace{10mm}

	\rule{8mm}{1.2pt}
\end{center}
% \thispagestyle{empty}
\clearpage


\begin{questions}

	\question[2]
	Betrachte die folgenden zwei Bilder. \cref{fig:fig1} zeigt ein Muster, erzeugt durch zahlreiche Kreise und \cref{fig:fig2} zeigt zufälliges Rauschen (white noise).

	\begin{figure}[htbp]
		\centering
		\begin{minipage}[b]{0.45\textwidth}
			\centering
			\includegraphics[width=\textwidth]{EF/Kolmogorov/Exams/circles.pdf}
			\caption{Kreise}
			\label{fig:fig1}
		\end{minipage}
		\hspace{0.05\textwidth}
		\begin{minipage}[b]{0.45\textwidth}
			\centering
			\includegraphics[width=\textwidth]{EF/Kolmogorov/Exams/white_noise.png}
			\caption{Rauschen}
			\label{fig:fig2}
		\end{minipage}
	\end{figure}

	Welches dieser Bilder hat vermutlich eine höhere Komplexität im Sinne von Kolmogorov? Begründe deine Antwort.

\begin{solution}
	siehe Skript
\end{solution}

	\question[1] Bestimme die Länge der kürzesten binären Darstellung der Zahl $92340283$.

	\begin{solution}
		Die gesuchte Länge ist gegeben durch $\ceil{\log_2\lr{92340283 + 1}} = 27$.
	\end{solution}

	\question
	\begin{parts}
		\part[1]
		Formuliere die Definition der plain Kolmogorov-Komplexität $C(w)$ eines binären Wortes $w$.
		\begin{solution}
	Betrachte die Definition der Kolmogorov-Komplexität im Skript.
\end{solution}
		\part[1] Erkläre, warum die Wahl der Programmiersprache in der Definition der Kolmogorov-Komplexität nicht wesentlich ist. Erkläre auch, was in diesem Zusammenhang mit \enquote{nicht wesentlich} gemeint ist.
				\begin{solution}
	Betrachte den Satz zur Sprachinvarianz der Kolmogorov-Komplexität im Skript. Mit \enquote{nicht wesentlich} ist gemeint, dass sich die Formulierungen in verschiedenen Programmiersprachen lediglich um eine Konstante unterscheiden würden.
\end{solution}
		\droptotalpoints
	\end{parts}

	\question[2] Wir betrachten die Wörter $w_n = 0^n1^{n^2}$ für alle $n\in\N$. Welche der folgenden drei Programmschemata sind gültige \texttt{C'++}-Programme, und welche davon liefern asymptotisch die beste obere Schranke für die Kolmogorov-Komplexität von $w_n$? Begründe deine Antwort.

	\textbf{Bemerkung:} Die Ausdrücke \verb|n| und \verb|n**2| zwischen den Escape-Sequenzen stehen für die konkreten Binärdarstellungen von $n$ beziehungsweise $n^2$. Das Präfix \verb|0b| legt fest, dass der nachfolgende Block als Binärdarstellung einer natürlichen Zahl interpretiert wird. Beachte bei der Beurteilung von \cref{listing:Cn}, dass ein gültiges \texttt{C'++}-Programm zwar denselben binären Block mehrfach, aber nicht zwei verschiedene binäre Blöcke enthalten darf.
\begin{lstlisting}[language=C++,caption=$A_n$]
    #include <iostream>
    #include <string>
    
    void An() {
        int x = 0b***n+++;
        int y = 0;
        for (int i = 1; i <= x; ++i) {
            std::cout << 0;
            y = y + x;
        }
        for (int j = 1; j <= y; ++j) {
            std::cout << 1;
        }
    }
\end{lstlisting}

\begin{lstlisting}[language=C++,caption=$B_n$]
    #include <iostream>
    #include <string>
    
    void Bn() {
        for (int i = 1; i <= 0b***n+++; ++i) {
            std::cout << 0;
        }
        for (int i = 1; i <= 0b***n+++; ++i) {
            for (int j = 1; j <= 0b***n+++; ++j) {
                std::cout << 1;
            }
        }
    }
\end{lstlisting}

\begin{lstlisting}[language=C++,caption=$C_n$,label=listing:Cn]
    #include <iostream>
    #include <string>
    
    void Cn() {
        for (int i = 1; i <= 0b***n+++; ++i) {
            std::cout << 0;
        }
        for (int i = 1; i <= 0b***n**2+++; ++i) { // n**2 steht im Schema für n*n
            std::cout << 1;
        }
    }
\end{lstlisting}
	\begin{solution}
		Die Programme $A_n$ und $B_n$ sind gültig und generieren jeweils $w_n$. In $B_n$ bezeichnen alle drei Vorkommen von \cppinline{0b***n+++} denselben binären Block und interpretieren ihn als die Zahl $n$; \texttt{g'++} speichert den Block nur einmal. Deshalb liefern beide Programme eine asymptotisch gleich gute obere Schranke:
		\begin{align*}
			C(w_n) &\leq \ceil{\log_2(n+1)} + c_A,\\
			C(w_n) &\leq \ceil{\log_2(n+1)} + c_B.
		\end{align*}
		Das Schema $C_n$ enthält dagegen die beiden verschiedenen Blöcke $bin(n)$ und $bin(n^2)$ und ist daher gemäss der Definition von \texttt{C'++} semantisch ungültig.
	\end{solution}


	\question[2] Sei $w_n := 0^{2^{2^{2n^4}}}$ für jedes $n\in\N, n>0$ ein binäres Wort. Bitte beachte unbedingt, dass Potenzausdrücke der Form $b^{p^q}$ gelesen werden müssen als
	\begin{align*}
		b^{p^q} := b^{\lr{p^q}}, \quad \text{also zum Beispiel} \quad 2^{3^2} := 2^{\lr{3^2}} = 2^9 = 512 \neq \lr{2^3}^2 = 8^2 = 64.
	\end{align*}
	Somit gilt
	\begin{align*}
		w_n := 0^{2^{2^{2n^4}}} := 0^{2^{\lr{2^{\lr{2n^4}}}}}.
	\end{align*}
	Gib eine (asymptotisch) möglichst gute obere Schranke für die Kolmogorov-Komplexität $C(w_n)$ von $w_n$ an, gemessen in der Länge $\abs{w_n}$ von $w_n$.
	\begin{solution}
		Das folgende \texttt{C'++} Programm erzeugt das Wort $w_n$:
		\begin{lstlisting}[language=C++,caption=$W_n$]
    #include <iostream>
    #include <cmath> // für die pow Funktion
    
    void Wn() {
		Natural k = 0b***n+++;
		k = 2 * pow(k, 4);
		k = pow(2, pow(2, k));
        for (int i = 1; i <= k; ++i) {
            std::cout << 0;
        }
    }
\end{lstlisting}
		Die Maschinencodes dieser Programme $W_n$ sind alle identisch und unabhängig von $w_n$ (und haben somit eine konstante binäre Länge) mit Ausnahme der Darstellung von $n$ in der Zeile Nummer 2. Somit ist die binäre Länge eines Programms $W_n$ gegeben durch
		\begin{align*}
			\ceil{\log_2(n+1)} + c
		\end{align*}
		für eine Konstante $c$. Damit lässt sich die Kolmogorov-Komplexität von $w_n$ abschätzen durch
		\begin{align*}
			C(w_n) \leq \ceil{\log_2(n+1)} + c \leq \blue{\log_2(n)} + c'
		\end{align*}
		für eine Konstante $c' := c + 1$. Die (binäre) Länge von $w_n$ ist
		\begin{align*}
			\abs{w_n} = 2^{2^{2n^4}}.
		\end{align*}
		Nach $n$ auflösen liefert
		\begin{align*}
			n = \lr{\frac{\log_2\lr{\log_2\lr{\abs{w_n}}}}{2}}^{1/4}
		\end{align*}
		und somit
		\begin{align*}
			\blue{\log_2(n)} & = \log_2\lr{\lr{\frac{\log_2\lr{\log_2\lr{\abs{w_n}}}}{2}}^{1/4}} = \\
			                 & = \frac{1}{4}\lr{\log_2\lr{\log_2\lr{\log_2\lr{\abs{w_n}}}}-1}.
		\end{align*}
		Dies ergibt die obere Schranke
		\begin{align*}
			C(w_n) \leq \frac{1}{4}\log_2\lr{\log_2\lr{\log_2\lr{\abs{w_n}}}} + c''
		\end{align*}
		für die Kolmogorov-Komplexität von $w_n$, wobei $c'':=c'-1/4$ eine Konstante ist.
	\end{solution}

	\question[1] Zeige, dass es für jede natürliche Zahl $n$ und jedes natürliche $i$ mit $i<n$ in dem Intervall $\ic{2^n}{2^{n+1} - 1}$ mindestens $2^n - 2^{n-i}$ natürliche Zahlen $x$ gibt mit $C(x) \geq n-i$.
	\begin{solution}
		Das Intervall $\ic{2^n}{2^{n+1} - 1}$ enthält genau $2^n$ natürliche Zahlen mit einer kürzesten binären Darstellungslänge von $n + 1$.
		Die Anzahl aller Maschinencodes der Länge kleiner als $n-i$ ist
		\begin{align*}
			\sum_{k=0}^{n-i-1}2^k=2^{n-i}-1.
		\end{align*}
		Jeder Maschinencode generiert höchstens eine Binärdarstellung und damit höchstens eine natürliche Zahl. Somit haben höchstens $2^{n-i}-1$ Zahlen im betrachteten Intervall eine Kolmogorov-Komplexität kleiner als $n-i$. Für die Anzahl $m$ der übrigen Zahlen gilt daher
		\begin{align*}
			m\geq2^n-(2^{n-i}-1)>2^n-2^{n-i}.
		\end{align*}
	\end{solution}
	
	\question[2] Beweise die folgende Behauptung:
	\begin{center}
		Für jede natürliche Zahl $n>0$ existiert mindestens ein binäres Wort $w$ der Länge $n$, für das gilt
		\begin{align*}
			C(w)\geq n = \abs{w},
		\end{align*}
		das heisst, für jede Länge $n$ existiert mindestens ein Wort dieser Länge, das sich nicht komprimieren lässt.
	\end{center}
	\begin{solution}
		Siehe das Theorem zur Existenz nichtkomprimierbarer Wörter im Skript.
	\end{solution}
\end{questions}
